\chapter{ABSTRAK}

\begin{singlespace}
\noindent
Navigasi autonomi luar jalan memberikan cabaran yang ketara disebabkan oleh rupa bumi yang tidak berstruktur, pelbagai ketinggian, dan tumbuh-tumbuhan yang padat. Algoritma perancangan laluan dua dimensi tradisional sering menghadapi kesukaran untuk memastikan keselamatan platform dan kecekapan pergerakan dalam persekitaran tiga dimensi yang kompleks ini. Tesis ini mencadangkan rangka kerja perancangan laluan peka-rupa bumi (\textit{terrain-aware}) baharu yang mengintegrasikan perwakilan topologi 2.5D setempat dengan penapisan lengkung licin bagi membolehkan navigasi yang selamat dan optimum. Metodologi teras memperkenalkan algoritma carian A* yang diubah suai dan menampilkan formulasi Kos Lintas Langkah (STC) pelbagai varian. STC ini sangat membebankan cerun yang curam---yang boleh dilaraskan melalui parameter pemberat cerun ($\lambda$)---serta gangguan tumbuh-tumbuhan, seterusnya menghalakan kenderaan dengan berkesan di sepanjang kontur yang selamat secara matematik. Untuk merapatkan jurang antara carian diskret berasaskan grid dengan pelaksanaan kenderaan yang dinamik, lapisan penapisan B-spline khusus diperkenalkan. Lapisan pascapemprosesan ini melicinkan trajektori secara iteratif sambil mengambil sampel kecerunan yang mendasari secara ketat untuk menjamin bahawa lengkung akhir tidak melanggar ambang keupayaan mendaki fizikal kenderaan. Simulasi menyeluruh merentas peta sintetik yang diseragamkan (rupa bumi kecerunan rendah, pelepasan sederhana dan pelepasan tinggi) menunjukkan keteguhan rangka kerja ini. Keputusan kuantitatif menunjukkan bahawa penapisan B-spline mengurangkan jumlah panjang laluan ruang sebanyak 5.06--6.28\%, proksi masa perjalanan sebanyak 4.98--5.40\%, dan purata sudut pusingan sebanyak 87.98--89.03\% berbanding dengan output A* mentah. Menariknya, pengoptimuman ini dicapai dengan sifar pelanggaran cerun dan menanggung kurang daripada 13 milisaat tambahan overhed pengiraan, sekali gus mengesahkan kesesuaiannya untuk penggunaan lapangan masa nyata. Rangka kerja yang dicadangkan ini memajukan secara ketara navigasi autonomi luar jalan dengan mengimbangi secara eksplisit kekangan keselamatan topologi dengan keberkesanan laluan serta kebolehlaksanaan dinamik yang berterusan.
\end{singlespace}